% =============================================================================
% PREÁMBULO — según resolución formal de Trabajo Final UCASAL
% =============================================================================

% ---- Codificación y lenguaje -------------------------------------------------
\usepackage[utf8]{inputenc}
\usepackage[T1]{fontenc}
\usepackage[spanish, es-tabla, es-ucroman]{babel}

% ---- Tipografía: Times New Roman 12 pt (especificado en resolución) ----------
\usepackage{newtxtext, newtxmath}

% ---- Página: A4, todos los márgenes 2.5 cm ----------------------------------
\usepackage[a4paper,
            top=2.5cm, bottom=2.5cm,
            left=2.5cm, right=2.5cm]{geometry}

% ---- Interlineado 1.5 + espaciado superior de 6 pt --------------------------
\usepackage{setspace}
\onehalfspacing
\setlength{\parindent}{6mm}
\setlength{\parskip}{6pt}

% ---- Alineación justificada (default LaTeX) ---------------------------------
% No se requiere paquete extra; LaTeX justifica por defecto.

% ---- Notas al pie: TNR 10 pt, interlineado simple, sin espaciado extra ------
\usepackage[bottom]{footmisc}
\makeatletter
\renewcommand\footnotesize{\@setfontsize\footnotesize{10pt}{12pt}}
\makeatother
\setlength{\footnotesep}{0pt}
\renewcommand{\footnotelayout}{\setstretch{1}}

% ---- Colores -----------------------------------------------------------------
\usepackage[table, xcdraw]{xcolor}
\definecolor{Micolor1}{RGB}{130,30,15}   % color institucional UCASAL
\definecolor{codebg}{rgb}{0.95,0.95,0.95}
\definecolor{comment}{rgb}{0.15,0.79,0.13}
\definecolor{keyword}{rgb}{0,0,1}
\definecolor{string}{rgb}{0.6,0.1,0.1}

% ---- Figuras -----------------------------------------------------------------
\usepackage{graphicx}
\graphicspath{{figures/}{assets/}{overleaft/}}
\usepackage{float}
\usepackage{subfigure}

% ---- Captions: TNR 10 pt, centrado, interlineado simple ---------------------
%   Tablas  → título ARRIBA  (position=top)
%   Figuras → título ABAJO   (position=bottom, default)
\usepackage{caption}
\DeclareCaptionFont{capten}{\fontsize{10pt}{12pt}\selectfont}
\captionsetup{
    font       = capten,
    labelfont  = {capten, bf},
    labelsep   = period,
    justification = centering,
    singlelinecheck = false,
    aboveskip  = 4pt,
    belowskip  = 0pt
}
\captionsetup[table]{position=top,  aboveskip=0pt, belowskip=4pt}
\captionsetup[figure]{position=bottom}

% ---- Matemáticas -------------------------------------------------------------
\usepackage{mathtools}
\usepackage{amsmath}

% ---- Tablas ------------------------------------------------------------------
\usepackage{booktabs}
\usepackage{longtable}
\usepackage{multirow}
\usepackage{array}
\usepackage{tabularx}
\usepackage{multicol}

% ---- TikZ y pgfplots ---------------------------------------------------------
\usepackage{tikz}
\usepackage{pgfplots}
\pgfplotsset{compat=1.18}
\usepackage{pgfplotstable}
\usetikzlibrary{trees, calc, positioning, chains, arrows.meta,
                shapes.misc, backgrounds, shapes.geometric}

% ---- Código fuente -----------------------------------------------------------
\usepackage{listings}
\lstset{
    backgroundcolor=\color{codebg},
    basicstyle=\ttfamily\footnotesize,
    commentstyle=\color{comment},
    keywordstyle=\color{keyword},
    stringstyle=\color{string},
    frame=single,
    breaklines=true,
    showstringspaces=false,
    numbers=left,
    numberstyle=\tiny\color{gray},
    captionpos=b,
    language=Python,
    morekeywords={self}
}

% ---- Misceláneos -------------------------------------------------------------
\usepackage{calc}
\usepackage{enumerate}
\usepackage{enumitem}
\usepackage{indentfirst}
\usepackage{siunitx}
\usepackage{url}
\usepackage{lastpage}
\usepackage{changepage}
\usepackage{pdfpages}
\usepackage{emptypage}
\setcounter{tocdepth}{2}      % máx. 3 niveles: capítulo, sección, subsección

% ---- Bibliografía: normas APA última edición (biblatex + biber) -------------
\usepackage{csquotes}
\usepackage[
    backend  = biber,
    style    = apa,
    sortlocale = es_ES,
    natbib   = true,      % habilita \citep, \citet compatibles con el texto
    backref  = true       % muestra "Citado en p. X" en la bibliografía
]{biblatex}
\addbibresource{bibliography.bib}

% Reemplaza "vid. pág." por "Citado en p./pp."
\DefineBibliographyStrings{spanish}{%
    backrefpage  = {Citado en p\adddot},
    backrefpages = {Citado en pp\adddot},
}

% ---- Hipervínculos (en negro, sin bordes) -----------------------------------
\usepackage{hyperref}
\hypersetup{
    colorlinks = true,
    linkcolor  = black,
    citecolor  = Micolor1,   % citas → color institucional UCASAL (burdeos)
    urlcolor   = black,
    pdfborder  = {0 0 0}
}
% pdftitle y pdfauthor se establecen al inicio del documento,
% cuando los comandos \tituloTesis y \autorTesis ya están definidos.
\AtBeginDocument{
    \hypersetup{pdftitle={\tituloTesis}, pdfauthor={\autorTesis}}
}

% ---- Encabezado y pie de página --------------------------------------------
%   Encabezado: Nombre de la Tesis (izq.) | Autor (der.)
%   Pie:        Número de página centrado
\usepackage{fancyhdr}
\pagestyle{fancy}
\fancyhf{}
\renewcommand{\chaptermark}[1]{\markboth{\chaptername~\thechapter.\enspace #1}{}}
\lhead{\small\nouppercase{\leftmark}}
\lfoot{\small\autorTesis}
\rfoot{\small\thepage}
\renewcommand{\headrulewidth}{1pt}
\renewcommand{\footrulewidth}{1pt}
\renewcommand{\headrule}{{\color{Micolor1}\hbox to\headwidth{\leaders\hrule height\headrulewidth\hfill}}}
\renewcommand{\footrule}{{\color{Micolor1}\hbox to\headwidth{\leaders\hrule height\footrulewidth\hfill}}}
\setlength{\headheight}{15pt}
% Estilo plain (primera pagina de capitulo): pie con regla, sin encabezado
\fancypagestyle{plain}{%
  \fancyhf{}%
  \lfoot{\small\autorTesis}%
  \rfoot{\small\thepage}%
  \renewcommand{\headrulewidth}{0pt}%
  \renewcommand{\footrulewidth}{1pt}%
  \renewcommand{\footrule}{{\color{Micolor1}\hbox to\headwidth{\leaders\hrule height\footrulewidth\hfill}}}%
}

% ---- Estilo de capítulos (nuevo en página nueva, con color UCASAL) ----------
\usepackage{titlesec}
% Nivel 1 — Capítulo
% Nivel 1 -- Capitulo: etiqueta "Capitulo N" + titulo grande + regla UCASAL
\titleformat{\chapter}[display]
  {\normalfont\color{Micolor1}}
  {\large\bfseries\chaptername\ \thechapter}
  {6pt}
  {\Huge\bfseries}
  [\vspace{6pt}{\color{Micolor1}\titlerule[1.5pt]}]
\titlespacing*{\chapter}{0pt}{0pt}{20pt}
% Nivel 2 -- Seccion con color UCASAL
\titleformat{\section}
  {\normalfont\large\bfseries\color{Micolor1}}
  {\thesection}{0.8em}{}
% Nivel 3 -- Subseccion
\titleformat{\subsection}
  {\normalfont\normalsize\bfseries}
  {\thesubsection}{0.8em}{}

% Nivel 2 — Sección
\titlespacing*{\section}{0pt}{12pt}{6pt}

% Nivel 3 — Subsección (último nivel permitido)
\titlespacing*{\subsection}{0pt}{10pt}{4pt}
